\documentclass[11pt]{article}
\usepackage[a4paper,margin=25mm]{geometry}
\usepackage{amsmath,amssymb,amsthm,bm}
\usepackage{booktabs}
\usepackage{xcolor}
\usepackage{microtype}
\usepackage[hidelinks]{hyperref}

\newcommand{\E}{\mathbb E}
\newcommand{\Var}{\operatorname{Var}}
\newcommand{\Sthr}{S_{\rm thr}}
\newtheorem{proposition}{Proposition}
\newtheorem{corollary}{Corollary}
\theoremstyle{remark}
\newtheorem{remark}{Remark}

\title{Firm Life Cycles and the Size--Variance Relation\\
\large The relative GLV under the unbalanced panel of Moran, Secchi and Bouchaud}
\author{}
\date{2 October 2026}

\begin{document}
\maketitle

\begin{center}
\fcolorbox{black!35}{blue!4}{\begin{minipage}{0.91\textwidth}
\textbf{Main results.}
\begin{enumerate}
  \item \emph{Protocol.} The empirical estimator is an unbalanced panel: every firm contributes over its own
        listed life.  The thesis estimator keeps only firms listed for a whole window, a set that empties as
        the window grows, so its $\beta$ has no long-window limit.  Counting each listed spell of a firm as one
        firm removes the window dependence, and only spells cut by the window edges depend on it.
  \item \emph{Single-firm reduction.} A firm driven by a Gaussian fitness with its measured mean $m_i$ and
        amplitude $f_i$, and no network, reproduces $\beta$ and $\zeta_q/\zeta_1$ of the full model under both
        protocols, and the slow window convergence.  The network enters only through the law of $(m,f)$ and
        the correlation time $\tau$.
  \item \emph{The $(m,f)$ law.} Mean-field theory gives $m\mid\kappa$ Gaussian with mean $1-\bar g-\mu\kappa\bar B$
        and variance $\sigma^2\kappa Q_T$, and $f=\sigma\sqrt{\kappa(Q_0-Q_T)}$: five order parameters and the degree
        law.  It holds firm by firm in simulation, and with the order parameters from the HDMFT the chain from
        $(\mu,\sigma,\lambda)$ and the degree law to $\beta$ and $\zeta_q/\zeta_1$ is closed with no simulation
        input ($\alpha=6$: $\beta=0.403$ against $0.422$, ratios within $0.03$).
  \item \emph{Lifetime identity.} For a spell that enters and leaves at the listing threshold, exactly,
        $\sigma=\sqrt{\pi/2}\,\bigl[2\ln(S_{\max}/\Sthr)+V_p\bigr]/n$.  The rise and fall carry $74$--$100\%$
        of the volatility on the decline branch, so $\sigma$ is essentially (log range)/(lifetime).
  \item \emph{Exponents are lifetime exponents.} $\zeta_q\simeq\zeta_q[n^{-1}]-q/\ell$ with
        $\ell=\ln(S_{\max}/\Sthr)$, verified to $0.02$.  The ratio curve bends because the dispersion of
        lifetimes grows with size: small firms live through regular fitness excursions, large firms exit only on
        rare deep dips.  This generalises the degree criterion of the $\beta$ note.
  \item \emph{Lifetime law.} Rice's formula for the fitness zero crossings gives the lifetime up to a merging
        factor $\approx2$; lifetime grows like $\exp(S^2/2f^2)$, so $\beta$ is an effective exponent across a
        crossover that moves as $\sqrt{\ln T}$.
  \item \emph{Degree channel.} On heavy-tailed graphs the flattening of the ratio curve is caused by the
        heterogeneity of fitness amplitudes $f\propto\sqrt\kappa$; making $f$ homogeneous restores the plain
        life-cycle bend.
\end{enumerate}
\end{minipage}}
\end{center}

\section{Two panels}

Moran, Secchi and Bouchaud (2024, Sec.~3.1) use Compustat 1961--2020.  Every firm with at least two annual
growth rates is kept (24{,}233 firms, between $2$ and $236$ growth rates each, $46$ on average).  Its volatility
is $\sigma_i=\sqrt{\pi/2}\,T_i^{-1}\sum_t|g_{it}-\bar g_i|$ over its own $T_i$ observations, its size the
lifetime average $\bar S_i$, and firms are binned by $\bar S_i$.  The higher exponents come from
$\E[\sigma^q\mid\bar S]\sim\bar S^{-\zeta_q}$.  A robustness sample keeps firms with at least $20$ growth rates.

The thesis estimator (\texttt{msb.size\_volatility}) is the same except for one step: it keeps only firms whose
share stays above $10^{-6}$ throughout the window $[180,240]$, about $13$--$18\%$ of the firms.  We call this
the \emph{survivor} (balanced) panel.  Its survivors are firms that did not dip over $T$, and their number
vanishes as $T\to\infty$, so the survivor $\beta$ has no long-window limit.  At $\alpha=2.5$ it drifts from
$0.414$ to $0.354$ between $T=60$ and $480$.

In the model a firm whose fitness turns negative decays to the immigration floor and climbs back when it turns
positive.  Over a few hundred time units $80$--$95\%$ of firms do this at least once, so a node behaves as a
niche that successive firms occupy, and the immigration term acts as entry.  We define a \emph{spell} as a
maximal stretch with $S>\Sthr=8\lambda$ (the share $10^{-6}$ at $\lambda=10^{-3}$, $N=8000$) and treat each spell
as one firm: this is the \emph{unbalanced} panel, implemented as \texttt{msb.spell\_volatility}.  Only spells cut
by the window edges depend on the window.  In long runs the spell $\beta$ converges: at $\alpha=2.2$ it is
$0.110, 0.081, 0.084, 0.080$ for $T=240,480,960,1920$, and at $\alpha=6$ it is $0.461, 0.422, 0.391, 0.374$,
with the cut-off share falling to $3$--$4\%$.  Section~\ref{sec:rice} explains why the approach is slow.

\section{Single-firm reduction}
\label{sec:reduction}

At zero reciprocity the effective process of a firm is (notation of the $\beta$ note)
\begin{equation}
  \frac{d\ln S}{dt}=\Phi(t)-S+\lambda\bigl(S^{-1}-1\bigr),\qquad
  \Phi_i(t)=1-g(t)-(\alpha S(t))_i ,
  \label{eq:process}
\end{equation}
and $\Phi_i$ can be read off the simulation exactly at every output time.

\paragraph{Fitness statistics.}  Table~\ref{tab:fitness} summarises them.  Within a firm, $\Phi_i=m_i+f_i\eta_i(t)$
with $\eta_i$ of unit variance.  Its autocorrelation is smooth at short lags and well described by
$\rho(L)=(1+L/\tau)e^{-L/\tau}$.  For listed firms with $S>1$ the time-mean size equals the time-mean fitness
while listed (ratio $1.001$ at $\alpha=6$, $1.003$ at $\alpha=3.5$), the size-is-mean-fitness law of the $\beta$
note.

\begin{table}[h]
\centering\small\setlength{\tabcolsep}{4pt}
\begin{tabular}{lcccccccc}
\toprule
graph, $\sigma$ & $\bar m$ & sd $m$ & skew $m$ & exkurt $m$ & $\bar f\pm$ sd & corr$(m,f)$ & ACF $1/e$ & $\tau$ fit\\
\midrule
$\alpha=6$, $1.75$   & $-0.95$ & $3.27$ & $-0.11$ & $0.46$  & $1.82\pm0.41$ & $-0.09$ & $7.5$ & $3.51$\\
$\alpha=3.5$, $1.9$  & $-1.73$ & $3.70$ & $-0.87$ & $5.4$   & $3.45\pm1.08$ & $-0.29$ & $3.5$ & $1.65$\\
$\alpha=2.5$, $1.75$ & $-0.97$ & $4.36$ & $-4.56$ & $56.9$  & $2.12\pm1.15$ & $-0.42$ & $4.0$ & $1.89$\\
\bottomrule
\end{tabular}
\caption{Exact fitness of the full model, $N=8000$, $\lambda=10^{-3}$, eight economies, window $[180,660]$.  On
heavy tails, hubs combine a large amplitude ($f\propto\sqrt\kappa$) with a very negative mean (their load
$\propto\kappa$), hence the skewed $m$ and the negative correlation.  Source:
\texttt{scripts/lifecycle/fitness\_stats.py}.}
\label{tab:fitness}
\end{table}

\paragraph{Surrogate.}  Replace the network by independent one-firm processes~\eqref{eq:process} with
$\Phi_i=m_i+f_i\eta_i$, where $(m_i,f_i)$ are the measured values and $\eta_i$ is a Gaussian process with the
fitted $\rho$ (an Ornstein--Uhlenbeck cascade).  The same estimators are then applied.

\begin{table}[h]
\centering\small\setlength{\tabcolsep}{4pt}
\begin{tabular}{ll|cc|cc}
\toprule
 & & \multicolumn{2}{c|}{full model} & \multicolumn{2}{c}{single-firm surrogate}\\
graph & panel & $\beta$ & $\zeta_q/\zeta_1$, $q=2,3,4$ & $\beta$ & $\zeta_q/\zeta_1$\\
\midrule
$\alpha=6$, $\sigma=1.75$ & survivor & $0.894$ & $1.96,\,2.87,\,3.70$ & $0.862$ & $1.95,\,2.82,\,3.62$\\
                          & spells   & $0.422$ & $1.86,\,2.55,\,3.07$ & $0.439$ & $1.79,\,2.39,\,2.82$\\
$\alpha=3.5$, $\sigma=1.9$ & spells  & $0.214$ & $1.80,\,2.31,\,2.54$ & $0.231$ & $1.79,\,2.35,\,2.69$\\
$\alpha=2.5$, $\sigma=1.75$ & survivor & $0.414$ & $1.87,\,2.58,\,3.15$ & $0.474$ & $1.87,\,2.60,\,3.21$\\
                          & spells   & $0.198$ & $1.49,\,1.48,\,1.14$ & $0.185$ & $1.52,\,1.67,\,1.56$\\
\midrule
\multicolumn{2}{l|}{$\alpha=6$, spells, $T=240,480,960,1920$} & \multicolumn{2}{c|}{$0.461,\,0.422,\,0.391,\,0.374$}
 & \multicolumn{2}{c}{$0.454,\,0.439,\,0.368,\,0.348$}\\
\bottomrule
\end{tabular}
\caption{The surrogate against the full model.  Data: $\beta\approx0.20$, ratios $2.0,\,2.6,\,2.9$.  Sources:
\texttt{scripts/lifecycle/surrogate.py}, \texttt{survivor\_check.py}, \texttt{window\_surrogate.py}.}
\label{tab:surrogate}
\end{table}

\begin{proposition}[Single-firm reduction]\label{prop:reduction}
Under both panels, $\beta$, the ratio curve and the window dependence are properties of the one-firm
process~\eqref{eq:process} driven by a Gaussian fitness.  The graph enters only through the joint law of $(m,f)$
and the correlation time.
\end{proposition}
This is the numerical content of Table~\ref{tab:surrogate}: agreement to $0.06$ in $\beta$ and to $0.25$ or
better in the ratios, except at $q=4$ on the flattest heavy-tail curve ($1.56$ against $1.14$).  It is what the cavity picture predicts at leading order
in $1/C$, where the field on a firm is Gaussian and independent of its own size.  Everything below is derived for
the one-firm process.

\section{The law of $(m,f)$ from mean-field theory}
\label{sec:law}

\begin{proposition}[Fitness law]\label{prop:law}
At $\gamma=0$ and leading order in $1/C$, the fitness of a firm of degree class $\kappa=k/C$ is
$\Phi(t)=1-g(t)-\mu\kappa\bar B+\sigma\sqrt\kappa\,\xi(t)$ with $\E[\xi(t)\xi(t')]=Q(t-t')$, where
$\bar B=\sum_jw_j\bar S_j$, $Q(L)=\sum_jw_j\overline{S_j(t)S_j(t+L)}$ and $w_j\propto k_j$.  Writing
$Q_T=\sum_jw_j\bar S_j^2$ for the part of $Q$ that is static over the window,
\begin{gather}
  m\mid\kappa\sim\mathcal N\bigl(r_\kappa,\;\sigma^2\kappa Q_T\bigr),\qquad r_\kappa=1-\bar g-\mu\kappa\bar B,
  \nonumber\\
  f(\kappa)=\sigma\sqrt{\kappa\,(Q_0-Q_T)},\qquad \rho(L)=\frac{Q(L)-Q_T}{Q_0-Q_T}.
  \label{eq:law}
\end{gather}
\end{proposition}
\begin{proof}
$(\alpha S)_i=(\mu/C)\sum_jA_{ij}S_j+(\sigma/\sqrt C)\sum_jA_{ij}z_{ij}S_j$.  In a configuration-model graph the
$\kappa C$ neighbours are drawn with probability $\propto k_j$, so the first sum is $\mu\kappa\bar B$ up to
$O(C^{-1/2})$.  At $\gamma=0$ the couplings $z_{ij}$ are independent of the neighbours' trajectories at leading
order, so the second sum is Gaussian with covariance $(\sigma^2/C)\sum_jA_{ij}S_j(t)S_j(t')\to\sigma^2\kappa Q(t,t')$.
Its window mean has variance $\sigma^2\kappa Q_T$ and the deviations from it have variance $\sigma^2\kappa(Q_0-Q_T)$.
\end{proof}

Five numbers ($\bar g$, $\bar B$, $Q_0$, $Q_T$ and the correlation time of $\rho$) and the degree law fix the
fitness of every firm.  Three consequences:
\begin{enumerate}
  \item \emph{Amplitude heterogeneity is degree heterogeneity.}  $f\propto\sqrt\kappa$ with no free parameter;
        this is the spread that flattens the heavy-tail ratio curve.
  \item \emph{Shape of $m$.}  The mean falls linearly in $\kappa$ (load) while the spread grows as $\sqrt\kappa$, so
        $\mathrm{corr}(m,f)<0$, and on a Pareto degree law the mixture over $\kappa$ skews $m$ to the left.
  \item \emph{Listing by degree.}  With total fitness variance $\sigma^2\kappa Q_0$, the fraction of time listed is
        $P(\text{listed}\mid\kappa)\simeq\Phi_{\mathcal N}\bigl(r_\kappa/\sigma\sqrt{\kappa Q_0}\bigr)$, which for hubs
        decays as a Gaussian in $\sqrt\kappa$: the survival thinning of the degree tail measured in the $\beta$
        note follows.
\end{enumerate}

\paragraph{Cavity check.}  With the order parameters measured as degree-weighted averages in the simulation, the
law predicts every firm's $(m_i,f_i)$ from its degree alone:

\begin{center}
\footnotesize\setlength{\tabcolsep}{3pt}
\begin{tabular}{l|cccc|c|ccc}
\toprule
 & \multicolumn{4}{c|}{$(m-r_\kappa)/\sigma\sqrt{\kappa Q_T}$} & $f/f(\kappa)$ & sd $m$ & skew $m$ & corr$(m,f)$\\
graph, $\sigma$ & mean & sd & skew & exkurt & median (IQR) & \multicolumn{3}{c}{measured / law}\\
\midrule
$\alpha=6$, $1.75$   & $-0.003$ & $1.004$ & $-0.01$ & $0.23$ & $0.97$ ($0.87$--$1.09$) & $3.27/3.23$ & $-0.11/-0.09$ & $-0.09/-0.12$\\
$\alpha=3.5$, $1.9$  & $0.006$  & $1.006$ & $-0.02$ & $0.28$ & $0.98$ ($0.88$--$1.09$) & $3.69/3.67$ & $-0.87/-0.87$ & $-0.29/-0.33$\\
$\alpha=2.5$, $1.75$ & $0.004$  & $0.997$ & $-0.03$ & $0.35$ & $0.96$ ($0.83$--$1.12$) & $4.36/4.25$ & $-4.56/-4.47$ & $-0.42/-0.41$\\
\bottomrule
\end{tabular}
\end{center}

The static field is a standard Gaussian once its predicted mean and spread are removed, the amplitude is
$\sigma\sqrt{\kappa(Q_0-Q_T)}$ to within $\pm13\%$, and the extreme skew of $m$ on the $\alpha=2.5$ graph ($-4.5$) is
reproduced by the mixture over degrees.  The listed fraction follows consequence~3 in its degree dependence (at
$\alpha=2.5$ it falls from $0.57$ to $0.30$ across the degree range, against $0.50$ to $0.31$) and sits
$0.04$--$0.07$ above it throughout, because a firm stays listed while it decays after its fitness turns negative,
the same delay that merges excursions in Section~\ref{sec:rice}.  Sources: \texttt{scripts/lifecycle/mf\_law.py},
\texttt{listed\_by\_degree.py}.

\paragraph{Closing the chain.}  Sampling $(m,f)$ from~\eqref{eq:law} and running the single-firm surrogate gives
the spell statistics at three levels of input:

\begin{center}
\footnotesize\setlength{\tabcolsep}{3pt}
\begin{tabular}{l|c|c|c|c}
\toprule
spells, $T=480$ & full model & measured $(m_i,f_i)$ & law (sim.\ order par.) & law (HDMFT)\\
\midrule
$\alpha=6$, $\sigma=1.75$  & $0.422$; $1.86,2.55,3.07$ & $0.439$; $1.79,2.39,2.82$ & $0.432$; $1.85,2.52,3.03$ & $0.403$; $1.86,2.56,3.10$\\
$\alpha=3.5$, $\sigma=1.9$ & $0.214$; $1.80,2.31,2.54$ & $0.231$; $1.79,2.35,2.69$ & $0.277$; $1.80,2.38,2.73$ & $0.285$; $1.81,2.43,2.86$\\
$\alpha=2.5$, $\sigma=1.75$ & $0.198$; $1.49,1.48,1.14$ & $0.185$; $1.52,1.67,1.56$ & $0.235$; $1.60,1.83,1.80$ & $0.216$; $1.63,1.91,1.92$\\
\bottomrule
\end{tabular}
\end{center}

Each entry is $\beta$; $\zeta_q/\zeta_1$ for $q=2,3,4$.  The last column has no simulation input: the HDMFT solved
on the degree sequence gives the order parameters, \eqref{eq:law} gives $(m,f)$, the one-firm process gives the
life cycles and the MSB estimator is applied to them.  At $\alpha=6$ it reproduces the full model to $0.02$ in
$\beta$ and $0.03$ in the ratios.  On every graph it agrees with the law built from the simulation's own order
parameters to $0.03$ in $\beta$.  The remaining differences have two sources, of which the first dominates.
\begin{itemize}
  \item \emph{Finite connectivity.}  At fixed degree the measured amplitudes spread by $\pm13\%$, because each
        firm's actual neighbourhood differs; the leading-order law has none of this spread.  Less amplitude spread
        means less flattening at $\alpha=2.5$ and a higher $\beta$ (by $0.06$ at $\alpha=3.5$), the same
        finite-connectivity correction as in the $\beta$ note.
  \item \emph{The closure itself.}  Solved on the realised degree sequence, the HDMFT order parameters give a law
        within about $10\%$ of the simulation's: load shifted by $0.01$--$0.12$, static spread $3$--$5\%$ smaller,
        amplitude between $-11\%$ and $+12\%$, and a field slightly rougher at $\alpha=6$ (correlation $0.74$ at
        lag $2$, against $0.85$).
\end{itemize}

\section{The lifetime identity}

\begin{proposition}[Lifetime identity]\label{prop:identity}
Let a spell start and end at $\Sthr$ with $n$ growth steps $g_t=\Delta\ln S$.  Then $\bar g=0$, so the MAD is
the total variation per step, and
\begin{equation}
  \sigma=\sqrt{\tfrac\pi2}\;\frac{2\ell+V_p}{n},\qquad \ell=\ln\frac{S_{\max}}{\Sthr},
  \label{eq:identity}
\end{equation}
where $V_p\ge0$ is the variation in excess of a single rise and fall.  For any firm, with entry size $S_a$ and
exit size $S_b$,
\begin{equation}
  \sigma\;\ge\;\sqrt{\tfrac\pi2}\;\frac{2\ln\bigl(S_{\max}/\max(S_a,S_b)\bigr)}{n}.
  \label{eq:bound}
\end{equation}
\end{proposition}
\begin{proof}
Split the steps at the peak into $A$ (a fraction $x$ of them) and $B$.  With $R_a=\ln(S_{\max}/S_a)$,
$R_b=\ln(S_{\max}/S_b)$ and $\bar g=(R_a-R_b)/n$,
$\sum_t|g_t-\bar g|\ge|\sum_A(g_t-\bar g)|+|\sum_B(g_t-\bar g)|=2[(1-x)R_a+xR_b]\ge2\min(R_a,R_b)$.  For a
complete spell $R_a=R_b=\ell$ and the bound is the rise-and-fall part of the total variation.
\end{proof}

The bound~\eqref{eq:bound} does not depend on the model and can be evaluated firm by firm in data.

In the surrogate at $\alpha=6$, \eqref{eq:identity} holds to a median relative error of $2.5\times10^{-3}$ (grid
discretisation) on the $87\%$ of spells that are complete at $T=480$.  The rise-and-fall term carries $100\%$,
$95\%$ and $74\%$ of $\sigma$ at the low end, middle and top of the decline branch.  Under the unbalanced panel,
then, the volatility of a firm is essentially its log range divided by its lifetime, and the plateau
fluctuations that the $\beta$ note analyses matter only for the largest firms.  The survivor panel sees the
opposite: by construction its firms never rise or fall, so it measures the plateau alone.

\section{Exponents as lifetime exponents}

When the rise and fall dominate, $\sigma\simeq c\,\ell/n$ with $c=2\sqrt{\pi/2}$.  If $\ell$ and $n$ are
independent at fixed $\bar S$,
\begin{equation}
  \zeta_q\simeq\zeta_q[n^{-1}]-\frac{d\ln\E[\ell^q\mid\bar S]}{d\ln\bar S}
  \simeq\zeta_q[n^{-1}]-\frac{q}{\ell},
  \label{eq:lifetime-rep}
\end{equation}
where $\zeta_q[n^{-1}]$ are the scaling exponents of the inverse lifetime and the last step uses
$S_{\max}\propto\bar S$.  So the size--variance exponent is a lifetime--size exponent, lowered by a logarithmic
correction set by the depth of entry.

\begin{table}[h]
\centering\footnotesize\setlength{\tabcolsep}{2.5pt}
\begin{tabular}{c|cc|c|c|c|c}
\toprule
$\lambda$ & $\ell$ & $\frac{d\ln\ell}{d\ln S}$ ($1/\ell$) & $\zeta_q$, transient & r.h.s.\ of \eqref{eq:lifetime-rep}
 & $\zeta_q[n^{-1}]$ & $\zeta_q/\zeta_1$ of $\sigma$\\
\midrule
$10^{-4}$ & $8.1$ & $0.107$ ($0.124$) & $.328,\,.546,\,.699,\,.815$ & $.329,\,.551,\,.712,\,.837$ & $.437,\,.766,\,1.03,\,1.27$ & $1.85,\,2.55,\,3.09$\\
$10^{-3}$ & $5.8$ & $0.150$ ($0.172$) & $.427,\,.718,\,.931,\,1.10$ & $.428,\,.722,\,.940,\,1.11$ & $.578,\,1.02,\,1.39,\,1.72$ & $1.88,\,2.61,\,3.20$\\
$10^{-2}$ & $3.5$ & $0.285$ ($0.287$) & $.518,\,.876,\,1.14,\,1.35$ & $.519,\,.877,\,1.14,\,1.36$ & $.803,\,1.45,\,2.00,\,2.49$ & $1.91,\,2.68,\,3.32$\\
\bottomrule
\end{tabular}
\caption{Complete spells on the decline branch, surrogate with the $\alpha=6$ fitness, $T=480$.
Eq.~\eqref{eq:lifetime-rep} holds to $0.02$.  At fixed fitness statistics $\lambda$ acts mainly through $\ell$ and
the lifetimes and its net effect on the ratios is small; in the full model $\lambda$ also changes $(m,f)$.
Source: \texttt{scripts/lifecycle/lifetime.py}.}
\label{tab:lifetime}
\end{table}

\section{Why the ratio curve bends}

The $\beta$ note showed that the degree channel bends $\zeta_q/\zeta_1$ below $q$ if and only if
$\Var(\ln\kappa\mid S)$ grows with size.  The argument uses only that the volatility is an amplitude times a random
factor, so it applies to any such factor.

\begin{corollary}[General multiscaling criterion]\label{cor:criterion}
If $\sigma=A(\bar S)\,X$ with $X$ random given $\bar S$, then
$\zeta_q-q\zeta_1=-\tfrac{q(q-1)}{2}\,\frac{d}{d\ln\bar S}\Var(\ln X\mid\bar S)+O(q^3)$, and the ratio curve
bends below $q$ if and only if the conditional log-variance of $X$ grows with size.  The degree channel has
$X=\sqrt\kappa$ and the life cycle $X=\ell/n$.
\end{corollary}

For the life cycle the relevant spread is that of $\ln n$.  At $\alpha=6$ (surrogate, $\lambda=10^{-3}$),
$\Var(\ln n\mid\bar S)$ rises from $0.20$ to $0.62$ along the decline branch while the mean lifetime grows from
$14$ to $81$.  The measured bend of the lifetime exponents, $\zeta_q[n^{-1}]-q\,\zeta_1[n^{-1}]=0,\,-0.13,\,-0.34,
\,-0.60$, has the predicted sign; the second-cumulant estimate ($0,\,-0.19,\,-0.58,\,-1.16$) is right at $q=2$ and
overshoots at higher $q$, as an expansion in $q^2$ must.

The mechanism is a crossover between two kinds of exit.  A small firm lives through a single positive excursion
of its fitness; such excursions have fairly regular lengths.  A large firm exits only on a rare deep dip, a
nearly memoryless event, so its lifetime is close to exponential, for which $\Var\ln n=\pi^2/6=1.64$.  The
spread therefore grows with size until the exits have become memoryless, and then it saturates: like the degree
channel, the life-cycle bend is a crossover and not an asymptotic law.

\section{The lifetime law}
\label{sec:rice}

For $\Phi=m+f\eta$ with $\rho''(0)=-1/\tau^2$, Rice's formula gives the rate of upcrossings of zero,
$\nu(u)=e^{-u^2/2}/(2\pi\tau)$ with $u=m/f$, and the mean excursion length $\bar L(u)=\Phi_{\mathcal
N}(u)/\nu(u)$.  A dip ends a spell only if it outlasts the decay time $\ln(S/\Sthr)/|\Phi|$, so neighbouring
excursions merge.

\begin{table}[h]
\centering\small\setlength{\tabcolsep}{4pt}
\begin{tabular}{c|cc|cc|cc|cc}
\toprule
$u$ & spells per time & Rice $\nu$ & mean $L$ & Rice $\bar L$ & time listed & $\Phi_{\mathcal N}(u)$ & $S$ while listed & $m$\\
\midrule
$-0.75$ & $0.021$ & $0.034$ & $12.4$ & $6.6$ & $0.26$ & $0.23$ & $0.60$ & $-1.34$\\
$0.25$  & $0.024$ & $0.044$ & $28.8$ & $13.6$ & $0.68$ & $0.60$ & $1.25$ & $0.43$\\
$0.75$  & $0.017$ & $0.034$ & $49.9$ & $22.6$ & $0.85$ & $0.77$ & $1.73$ & $1.29$\\
$1.25$  & $0.009$ & $0.021$ & $84.8$ & $43.1$ & $0.94$ & $0.89$ & $2.33$ & $2.14$\\
$1.75$  & $0.005$ & $0.010$ & $118.9$ & $98.0$ & $0.98$ & $0.96$ & $3.03$ & $2.96$\\
\bottomrule
\end{tabular}
\caption{Per firm, binned in $u=m/f$, surrogate with the $\alpha=6$ fitness, $T=480$.  Above $u\approx2$ the
lifetimes exceed the window.  Source: \texttt{scripts/lifecycle/rice.py}.}
\label{tab:rice}
\end{table}

Table~\ref{tab:rice} shows Rice's $u$-dependence with a merging factor: spells occur at $0.45$--$0.6\,\nu$ per
unit time and, for $u\le1.25$, last about twice as long as single excursions.  The listed fraction is close to $\Phi_{\mathcal N}(u)$ and,
for $u\gtrsim1$, the size while listed is $m$.  Two consequences follow.
\begin{enumerate}
  \item \emph{Gaussian lifetime law.} For $u\gg1$, $\ln\bar L\simeq u^2/2=\bar S^2/2f^2$: lifetime grows like a
        Gaussian in size, not as a power.  Through~\eqref{eq:lifetime-rep}, the unbalanced-panel $\beta$ is an
        effective exponent whose value depends on the fitted size range.
  \item \emph{Crossover and slow convergence.} Firms whose lifetime exceeds the panel length $T$ appear as one
        cut-off spell dominated by the plateau ($\sigma\propto1/S$); shorter-lived firms are dominated by the
        rise and fall.  The crossover size is $S_\times\simeq f\sqrt{2\ln(cT/2\pi\tau)}$ with $c\approx0.5$.  It
        moves only as $\sqrt{\ln T}$, which is why the spell $\beta$ converges so slowly in $T$
        (Table~\ref{tab:surrogate}, last row).
\end{enumerate}

\section{The degree channel is amplitude heterogeneity}

At $\alpha=2.5$ the spell ratio curve is flat.  Splitting $\sigma=\sqrt\kappa\,A\,\zeta$ as in the $\beta$ note:
the degree-normalised ratios are $1.82,\,2.38,\,2.68$, close to the data, while the degree moments grow faster
than $q$ ($e_q/e_1=1,\,2.2,\,3.5,\,4.7$) and lower $\zeta_1$ from $0.29$ to $0.18$.  Dividing the same absolute
bend by a smaller $\zeta_1$ flattens the curve.  The surrogate isolates the cause:

\begin{center}
\small
\begin{tabular}{lcc}
\toprule
$\alpha=2.5$, spells & $\beta$ & $\zeta_q/\zeta_1$\\
\midrule
full model & $0.198$ & $1.49,\,1.48,\,1.14$\\
surrogate, measured $(m_i,f_i)$ & $0.185$ & $1.52,\,1.67,\,1.56$\\
$f$ shuffled across firms & $0.116$ & $1.50,\,1.61,\,1.38$\\
$f$ homogeneous & $0.365$ & $1.90,\,2.67,\,3.29$\\
\bottomrule
\end{tabular}
\end{center}

Breaking the link between $m$ and $f$ keeps the curve flat; removing the spread of $f$ restores the plain
life-cycle bend.  The heavy-tail flattening is therefore the spread of fitness amplitudes, $f\propto\sqrt\kappa$:
at fixed size it adds lifetime dispersion on top of the life cycle, which is Corollary~\ref{cor:criterion} with
two random factors.  Source: \texttt{scripts/lifecycle/ablation.py}.

\section{Where the model meets the data}

A scan over $\alpha\in\{2.5,3,3.5,6\}$, $\sigma\in\{1.6,1.75,1.9\}$, $\lambda\in\{10^{-4},10^{-3},10^{-2}\}$
($\mu=1.76$, $N=8000$, eight economies, \texttt{scripts/scan\_spells.py}) shows the two targets pulling apart.
$\beta$ in the band needs heavy tails or large $\sigma$; a data-like ratio curve needs light tails, where the
amplitude spread is small.  They meet near $\alpha\approx3.5$, $\sigma\approx1.9$: with spells of at least $20$
growth rates and $\lambda=10^{-2}$, $\beta=0.19$ and $\zeta_q/\zeta_1=1.80,\,2.41,\,2.83$; with spells of at least
$2$ and $\lambda=10^{-3}$, $\beta=0.21$ and $1.80,\,2.31,\,2.54$.  On light tails raising $\sigma$ lowers $\beta$
without damaging the ratios, so the joint region probably extends past the scanned range.  Across all $36$
points the growth rates stay fat-tailed (excess kurtosis $3$--$16$).

\section{Interpretation, predictions and caveats}

\paragraph{Two panels, two mechanisms.}  In the model, the survivor panel measures plateau fluctuations: the
adiabatic law and degree channel of the $\beta$ note, with $\beta\to1$ on homogeneous graphs.  The unbalanced
panel measures life cycles: entry, growth, decline and exit, with $\beta$ a lifetime--size exponent and the
ratio curve set by the spread of lifetimes.  The mechanism behind the multiscaling depends on which firms are
counted, and the empirical protocol counts all of them.

\paragraph{Predictions for data.}
\begin{enumerate}
  \item \emph{Survivorship.} Restricting a panel to firms present throughout should raise $\beta$
        (in the model, $0.41$ against $0.20$ at $\alpha=2.5$).
  \item \emph{Lifetime bound.} Eq.~\eqref{eq:bound} gives a model-free lower bound on each firm's volatility
        from its life-cycle range and lifetime, which can be evaluated in Compustat.
  \item \emph{Age.} Removing the entry and exit years of each firm, or conditioning on age, should move
        $\zeta_q/\zeta_1$ towards $q$.
\end{enumerate}

\paragraph{Caveats.}
\begin{itemize}
  \item Model firms enter at $\Sthr\approx10^{-2}$ of the mean size and rise by $e^{\ell}\sim e^{5}$.  Real firms
        enter Compustat at their IPO size, so the model probably overstates the rise-and-fall share; how large it
        is in the data is the question prediction~3 answers.
  \item The listing threshold is part of the comparison, not of the model.  Here it follows the floor,
        $\Sthr=8\lambda$.
  \item The chain is closed at leading order in $1/C$.  The spread of amplitudes at fixed degree, which comes from
        quenched neighbourhoods, is not derived; it controls the extra flattening on heavy tails and lowers $\beta$ by
        up to $0.07$.  The merging factor of the Rice law is also not derived.
  \item All numbers are at one $\mu$, $N=8000$ and eight economies per point; the corrections in $N$ of the
        $\beta$ note have not been checked under the unbalanced panel.
\end{itemize}

\section*{Summary}
\begin{enumerate}
  \item The unbalanced panel has a long-window limit; the survivor panel does not.
  \item Both panels reduce to a single firm driven by a Gaussian fitness (Proposition~\ref{prop:reduction}).
  \item Under the unbalanced panel, volatility is log range over lifetime (Proposition~\ref{prop:identity}),
        and the exponents are lifetime exponents with a $q/\ell$ correction~\eqref{eq:lifetime-rep}.
  \item The ratio curve bends when the conditional log-variance of a random factor grows with size
        (Corollary~\ref{cor:criterion}); for the life cycle this is the crossover from regular excursions to
        memoryless exits, and on heavy tails the spread of amplitudes adds to it.
  \item Lifetime grows as $\exp(\bar S^2/2f^2)$, so $\beta$ is an effective exponent and its window
        convergence is logarithmic.
  \item The $(m,f)$ law follows from mean-field theory (Proposition~\ref{prop:law}); with HDMFT order parameters
        the chain to the MSB statistics has no simulation input.
  \item Open: the finite-connectivity spread of amplitudes, the merging factor, finite-$N$ corrections, and the
        empirical tests above.
\end{enumerate}

\end{document}
